\documentclass[UTF8,12pt,a4paper,fontset=fandol]{ctexart}

% 支持从项目根目录、深度学习目录或当前目录编译。
\makeatletter
\def\input@path{{../}{../../}}
\makeatother
\usepackage{researchnotes}

\title{知识蒸馏的原理与前置知识}
\author{}
\date{}

\begin{document}
\maketitle

\begin{abstract}
知识蒸馏通过模型产生的预测、表示或样本关系，为另一个模型提供训练信号。理解这一方法，既要掌握条件概率、交叉熵、梯度和经验风险，也要分清教师提供的信息、学生的表达能力以及蒸馏数据的覆盖范围。本文从一个三分类问题出发，逐步建立数学记号，推导经典输出蒸馏、温度缩放和高温近似，分析蒸馏可能改善学习的机制及其边界，再讨论特征、关系、在线、自监督、无原始数据、序列和生成模型中的蒸馏。全文将可证明的等式、用于理解的简化模型、算法设计和经验判断分开，并给出训练流程、数值检查、实验方案与练习解答，使读者能够独立阅读相关论文并实现可靠的基线。
\end{abstract}
\tableofcontents

\section{研究对象、基本问题与学习路线}
\label{sec:kd-introduction}

\subsection{从一个三分类例子理解蒸馏}
假设一个图像分类系统需要区分猫、狗和汽车。对某张猫的图片，人工标签只给出“猫”，可以写成向量 $(1,0,0)$。一个已经训练好的模型可能输出 $(0.70,0.25,0.05)$，另一个模型可能输出 $(0.70,0.05,0.25)$。二者对正确类别给出的概率相同，但对错误类别的判断不同：前者认为这张图片更像狗，后者认为它更像汽车。这种差别描述了模型在该输入附近形成的类别关系。人工标签没有直接给出这些关系，学生却可能从教师的完整输出中学习到它们。这里的概率首先是模型的估计，并不意味着真实世界中这张已经确定的图片有百分之二十五的概率变成狗。

\keyterm{知识蒸馏（knowledge distillation, KD）}是利用一个或多个模型提供的训练信号，训练另一个模型的技术集合。提供信号的模型称为\keyterm{教师模型（teacher model）}，接受信号的模型称为\keyterm{学生模型（student model）}。在经典压缩场景中，教师较大、推理较慢，学生较小、部署较便宜；但大小关系不是定义的一部分。相同结构的模型、同一模型在不同训练时刻的状态，甚至共同训练的多个模型，也可以承担教师和学生的角色。

蒸馏传递的对象通常是函数行为或表示结构，而不是把教师参数逐项复制到学生参数中。教师有一百层、学生只有十层时，两者的参数甚至没有逐项对应关系。只要对相同输入能定义可比较的输出，便可构造输出蒸馏；若能访问中间层，还可以比较特征或样本关系。经典温度软化方法的代表文献是 Hinton 等人的工作\cite{hinton2015}。本文使用自设数值例子和逐步推导解释这些目标，不把某篇论文在特定数据集上的收益当成普遍性能保证。

\subsection{蒸馏、压缩、迁移与伪标签}
\keyterm{模型压缩（model compression）}关心参数存储、计算量、内存或实际延迟的降低。剪枝删除连接或结构，量化降低数值精度，低秩分解改变矩阵表示，蒸馏则改变训练信号；这些方法可以组合，但并非同一种操作。蒸馏后仍部署与教师同样大小的模型，可能改善预测却没有压缩。学生参数量减少，也不必然带来同比例的速度提升，因为访存、算子实现、并行度和序列长度都会影响耗时。

\keyterm{迁移学习（transfer learning）}强调把已有任务或数据中获得的信息用于新的学习问题，可以通过参数初始化、冻结特征或重新训练实现。蒸馏是传递信息的一条路径，但可以发生在完全相同的任务上。\keyterm{伪标签（pseudo-label）}通常指模型为无标签样本产生的标签；若只保留教师概率最大的类别，就得到硬伪标签。保留整个概率分布则提供软监督。现代语言模型语境中的“蒸馏”还常指让教师生成文本，再以这些文本训练学生；这属于行为或序列层面的迁移，不能自动等同于对完整词元概率分布做 KL 匹配。

学习本文可沿四条相接的路线推进：先建立向量、概率、损失和反向传播的基础，再理解输出蒸馏及温度的数学作用，然后比较不同知识载体与训练组织方式，最后讨论语言模型、工程实现和实验判断。前置知识部分不要求读者已经熟悉信息论；用到的关键概念会在首次出现时定义。后文的算法和实验设计均是可执行方案，不表示已经运行并获得了相应结果。

\section{线性代数、微积分与计算图基础}
\label{sec:kd-mathematics}

\subsection{向量、矩阵和张量的维度}
设单个输入为列向量 $\mathbf x\in\mathbb R^d$，分类类别数为 $K\geq2$。一个线性分类层可以写成 $\mathbf z=\mathbf W\mathbf h+\mathbf b$，其中隐藏表示 $\mathbf h\in\mathbb R^m$、权重矩阵 $\mathbf W\in\mathbb R^{K\times m}$、偏置 $\mathbf b\in\mathbb R^K$，输出 $\mathbf z\in\mathbb R^K$。深度网络把多个线性或非线性变换复合起来，最终仍可输出这样一个长度为 $K$ 的向量。批量包含 $B$ 个样本时，若每行放置一个样本的输出，批量输出矩阵便属于 $\mathbb R^{B\times K}$。类别维度和批量维度在损失计算中承担不同作用，不能混为一次无差别平均。

图像中间特征常写成 $\mathbf F\in\mathbb R^{B\times C\times H\times W}$，四个维度依次为批量、通道、高度和宽度；语言模型隐藏状态常写成 $\mathbf H\in\mathbb R^{B\times L\times d_h}$，其中 $L$ 为词元位置数，$d_h$ 为隐藏维度。这样的多维数组称为\keyterm{张量（tensor）}。在本文的工程语境中，张量一词主要指带有明确轴含义的多维数值数组，不预设读者掌握微分几何中的张量定义。比较教师和学生特征之前，要先确定每一个轴是否具有可对应的含义。

向量的欧氏范数为 $\|\mathbf u\|_2=(\sum_j u_j^2)^{1/2}$，矩阵的 Frobenius 范数为 $\|\mathbf A\|_F=(\sum_{i,j}A_{ij}^2)^{1/2}$。平方误差蒸馏使用这些范数衡量数值差异；余弦相似度 $\mathbf u^\top\mathbf v/(\|\mathbf u\|_2\|\mathbf v\|_2)$ 则比较方向，要求两个向量非零。相同方向而不同长度的向量具有相同余弦相似度，但平方距离可能很大。因此，“表示相似”必须说明选择的是哪一种几何关系。

\subsection{梯度、链式法则与停止梯度}
设学生全部参数组成 $\theta\in\mathbb R^p$，训练目标为标量函数 $J(\theta)$。其梯度 $\nabla_\theta J$ 的第 $j$ 个分量是 $\partial J/\partial\theta_j$。梯度下降用
\begin{equation}
  \theta_{t+1}=\theta_t-\eta_t\nabla_\theta J(\theta_t)
  \label{eq:kd-gradient-descent}
\end{equation}
更新参数，其中 $t$ 是迭代索引，$\eta_t>0$ 是学习率。负梯度是当前点的局部下降方向；它并不保证在非凸网络目标上找到全局最优解。学习率过大可能越过下降区域，梯度很小也可能只是进入平坦区域，而不是已经得到最好的分类器。

如果 $J$ 通过输出 $\mathbf z_\theta$ 依赖参数，则链式法则给出
\begin{equation}
  \nabla_\theta J
  =\left(\frac{\partial\mathbf z_\theta}{\partial\theta}\right)^\top
    \nabla_{\mathbf z}J.
  \label{eq:kd-chain}
\end{equation}
矩阵 $\partial\mathbf z_\theta/\partial\theta\in\mathbb R^{K\times p}$ 称为 Jacobian 矩阵。\keyterm{反向传播（backpropagation）}按照计算图反向应用链式法则，高效计算这个乘积，一般不显式构造完整 Jacobian。蒸馏首先改变输出层或中间层收到的误差信号，再通过同样的链式法则改变学生参数；它不需要另外发明一套求导规则。

以教师输出 $\mathbf q_\phi$ 为目标时，经典离线蒸馏把教师参数 $\phi$ 固定。记 $\operatorname{sg}(\cdot)$ 为\keyterm{停止梯度（stop-gradient）}算子，它在前向计算中保持数值不变，在求导时把输入视为常量。损失 $\ell(\operatorname{sg}(\mathbf q_\phi),\mathbf p_\theta)$ 只更新学生。停止梯度、冻结参数以及把模型切换到推理模式是不同操作：前两者控制导数和参数更新，推理模式还可能控制随机失活与批量归一化统计量的使用，不能相互替代。

\subsection{二阶近似说明了什么}
对在邻域内二次连续可微的标量函数，参数扰动 $\Delta\theta$ 足够小时，有
\[
 J(\theta+\Delta\theta)=J(\theta)+\nabla J(\theta)^\top\Delta\theta
 +\frac12\Delta\theta^\top\nabla^2J(\theta)\Delta\theta
 +o(\|\Delta\theta\|_2^2).
\]
其中 $\nabla^2J$ 是 Hessian 矩阵，描述局部曲率。后文用类似的展开说明，高温概率匹配为何接近中心化 logit 的平方匹配。近似需要明确展开点和小量；“两个损失在某个极限下等价”不等于在任意温度、任意参数尺度下都可互换。这一区别会直接影响损失权重的选择。

\section{概率、统计与监督学习基础}
\label{sec:kd-probability}

\subsection{条件分布、样本与模型概率}
随机输入记为 $X$，随机标签记为 $Y\in\{1,\ldots,K\}$，其联合分布记为 $P$。小写 $x,y$ 表示一次实际取值。真实条件概率向量记为 $\mathbf r(x)$，其中 $r_k(x)=P(Y=k\mid X=x)$；对连续输入，可将其理解为一个几乎处处定义的正则条件分布。学生输出 $p_{\theta,k}(x)$，教师输出 $q_{\phi,k}(x)$，二者都是对真实条件规律的模型估计，通常不等于 $r_k(x)$。模型给出数值不意味着已经获得真实后验分布。

训练集写成 $\mathcal D_\ell=\{(x_i,y_i)\}_{i=1}^n$。基础统计分析通常假设这些样本独立同分布地来自 $P$，而真实时间序列、用户日志和重复采样未必满足这一假设。无标签输入集写成 $\mathcal D_u=\{u_j\}_{j=1}^m$，其输入分布记为 $Q_X$。蒸馏可在 $\mathcal D_\ell$、$\mathcal D_u$ 或两者的混合上进行；部署输入分布另记为 $P_X^{\mathrm{deploy}}$。区分这些分布，是理解覆盖不足和分布偏移的前提。

\keyterm{硬标签（hard label）}可以表示为独热向量 $\mathbf e_y$，其第 $y$ 个分量为一，其余为零。\keyterm{软标签（soft label）}是概率单纯形 $\Delta^{K-1}=\{\mathbf q:q_k\geq0,\sum_kq_k=1\}$ 中的一般向量。软标签不必来自教师，也可能来自多位标注者的比例或人为平滑。教师输出也不必是概率；特征蒸馏传递的是向量或张量，因此不能把所有蒸馏都理解成“把标签变软”。

\subsection{总体风险、经验风险与泛化}
对损失函数 $\ell$，\keyterm{总体风险（population risk）}与\keyterm{经验风险（empirical risk）}分别定义为
\begin{equation}
 R(\theta)=\mathbb E_{(X,Y)\sim P}\ell(p_\theta(X),Y),\qquad
 \widehat R_n(\theta)=\frac1n\sum_{i=1}^n\ell(p_\theta(x_i),y_i).
 \label{eq:kd-risk}
\end{equation}
训练直接计算的是经验风险，应用关心的是部署分布上的风险。训练误差低，只说明拟合了已观察到的样本，并不自动说明泛化良好。教师也会过拟合，所以“学生接近教师”与“学生接近真实规律”是两个不同命题。验证集用于选择温度、损失权重和停止时刻；测试集用于最后评估。反复根据测试结果修改蒸馏方案，会使测试集逐渐承担训练选择的角色。

\keyterm{小批量随机梯度下降（mini-batch stochastic gradient descent）}在每一步抽取样本子集，以其平均梯度近似经验风险梯度。在均匀抽样、当前参数固定且单样本梯度存在时，小批量平均梯度是全数据平均梯度的无偏估计。若按教师置信度不均匀采样而不校正权重，优化的通常已经是重加权风险。重加权可以有用，但要承认学习目标改变了，不能继续把结果解释为对原始样本分布的等权拟合。

\subsection{不确定性、置信度和校准}
最大类别概率 $\max_kq_k(x)$ 常被当作置信度指标，它只描述输出分布有多集中。\keyterm{概率校准（probability calibration）}要求在一组预测概率相近的事件中，实际发生比例与预测概率相符。一个模型总是用 $0.99$ 的概率预测，却只有约八成正确，便具有过度自信倾向。准确率比较的是最大概率类别是否正确，校准比较的是概率大小是否可信；两个教师准确率相同，也可能提供质量很不相同的软标签。

\keyterm{偶然不确定性（aleatoric uncertainty）}通常指观测噪声或任务本身存在的不可消除歧义，\keyterm{认知不确定性（epistemic uncertainty）}通常指有限数据和模型知识不足带来的不确定性。这是有用的建模区分，但单次 softmax 输出不能自动把二者分解出来。教师在模糊图片上给出分散概率，可能反映真实歧义；也可能只是模型没有学好。蒸馏损失本身无法鉴别这两种来源，需要标签、额外数据或有针对性的评估支持。

\section{信息论基础：熵、交叉熵与散度}
\label{sec:kd-information}

\subsection{熵和交叉熵的定义}
以下对数均为自然对数。对离散概率向量 $\mathbf q$，\keyterm{熵（entropy）}定义为 $H(\mathbf q)=-\sum_kq_k\log q_k$，约定 $0\log0=0$。熵衡量分布自身的不确定程度：确定分布的熵为零，$K$ 类均匀分布的熵为 $\log K$。熵高不意味着知识多或预测好；完全均匀的输出具有最高熵，却可能对当前分类问题没有区分能力。需要区分“分布不确定性”与“这个训练目标对学生是否有用”。

给定目标分布 $\mathbf q$ 和预测分布 $\mathbf p$，\keyterm{交叉熵（cross-entropy）}定义为
\begin{equation}
 H(\mathbf q,\mathbf p)=-\sum_{k=1}^Kq_k\log p_k.
 \label{eq:kd-crossentropy}
\end{equation}
如果目标是独热标签，便有 $H(\mathbf e_y,\mathbf p)=-\log p_y$。这也可以从最大似然得到：假设给定输入后标签按学生预测的类别分布产生，独立样本的条件似然是 $\prod_i p_{\theta,y_i}(x_i)$，取负对数后成为交叉熵之和。对固定样本数除以 $n$ 不改变最优点，但会改变梯度数值尺度。

\subsection{KL 散度及其非负性}
\begin{definition}[Kullback--Leibler 散度]
对两个离散概率向量，定义
\begin{equation}
 D_{\mathrm{KL}}(\mathbf q\|\mathbf p)
 =\sum_{k:q_k>0}q_k\log\frac{q_k}{p_k}.
 \label{eq:kd-kl}
\end{equation}
若存在 $q_k>0$ 而 $p_k=0$，则该值为正无穷；否则按上式求和。第一个参数决定加权和中的权重，交换参数一般会改变结果。
\end{definition}

\begin{proposition}[非负性和交叉熵分解]
当 $\mathbf p,\mathbf q$ 为同一有限类别集合上的概率分布时，$D_{\mathrm{KL}}(\mathbf q\|\mathbf p)\geq0$，等号成立当且仅当 $\mathbf p=\mathbf q$。当散度有限时，
\begin{equation}
 H(\mathbf q,\mathbf p)=H(\mathbf q)+D_{\mathrm{KL}}(\mathbf q\|\mathbf p).
 \label{eq:kd-ce-kl}
\end{equation}
\end{proposition}
\begin{proof}
散度无穷时非负性成立。否则，在 $S=\{k:q_k>0\}$ 上利用 $-\log u\geq1-u$，有
\[
 \sum_{k\in S}q_k\left(-\log\frac{p_k}{q_k}\right)
 \geq\sum_{k\in S}(q_k-p_k)
 =1-\sum_{k\in S}p_k\geq0.
\]
第一个不等式取等要求 $p_k=q_k$ 对所有 $k\in S$ 成立；这又使支持集外的 $p_k$ 全为零，故等号恰好对应两个分布相等。将 $\log(q_k/p_k)$ 拆成两个对数并分别求和，即得式\eqref{eq:kd-ce-kl}。
\end{proof}

固定教师时，$H(\mathbf q)$ 不依赖学生参数，因此最小化交叉熵与最小化正向 KL 对学生产生相同梯度；两种损失的数值本身并不相同。若教师也接受梯度更新，教师熵就不再是常量，上述优化等价性不能直接沿用。KL 也不是通常意义上的距离：它不对称，且一般不满足三角不等式。本文称 $D_{\mathrm{KL}}(\mathbf q\|\mathbf p)$ 为从教师到学生的正向 KL，称相反顺序为反向 KL；阅读论文时应以公式为准，避免名称约定不同造成误解。

\subsection{一个可手算的数值例子}
\begin{example}[同一正确率下的软分布差异]
设教师给出 $\mathbf q=(0.7,0.2,0.1)$，学生给出 $\mathbf p=(0.6,0.3,0.1)$，真实标签为第一类。硬标签交叉熵、软标签交叉熵和 KL 分别是多少？
\end{example}
\begin{solve}
硬标签交叉熵为 $-\log0.6\approx0.510826$。软标签交叉熵为
\[
 -0.7\log0.6-0.2\log0.3-0.1\log0.1\approx0.828631.
\]
教师熵为 $-0.7\log0.7-0.2\log0.2-0.1\log0.1\approx0.801819$，所以 KL 约为 $0.026812$。二者最大概率类别相同，但学生在第一类和第二类之间的概率分配仍不同。蒸馏可以在分类决策已经一致时继续产生梯度，因为其目标比类别一致更强。
\end{solve}

\section{神经网络分类、softmax 与监督梯度}
\label{sec:kd-softmax}

\subsection{从未归一化分数到概率}
分类网络在最后一层输出实数向量 $\mathbf z(x)$，其分量称为\keyterm{logit（未归一化类别分数）}。在二分类统计模型中，logit 常特指对数几率；在神经网络多分类语境中，则通常泛指 softmax 之前的分数。温度为 $\tau>0$ 时，定义
\begin{equation}
 \operatorname{softmax}_\tau(\mathbf z)_k
 =\frac{\exp(z_k/\tau)}{\sum_{j=1}^K\exp(z_j/\tau)}.
 \label{eq:kd-softmax}
\end{equation}
本文使用 $\tau$ 表示温度，以免与教师标记混淆。有限 logit 给出的每个概率严格为正，总和为一。softmax 并不会检验这些分数是否来自可靠模型，它只完成归一化。温度为一时省略温度上标。

对任意常数 $c$，有 $\operatorname{softmax}_\tau(\mathbf z+c\mathbf1)=\operatorname{softmax}_\tau(\mathbf z)$，其中 $\mathbf1$ 为全一向量。这是因为分子、分母都乘上了 $e^{c/\tau}$。因此概率只依赖相对分数，无法唯一确定所有 logit 的绝对值。实际计算时可减去最大分数，避免指数溢出；计算对数概率时应使用稳定的 log-sum-exp 形式
\[
 \log p_k=\frac{z_k}{\tau}-a-\log\sum_j\exp(z_j/\tau-a),
 \qquad a=\max_j z_j/\tau.
\]
这是精确代数变换，并非对目标函数的近似。

\subsection{softmax 的导数和交叉熵梯度}
直接对式\eqref{eq:kd-softmax}求导，利用商法则，得到
\begin{equation}
 \frac{\partial p_k}{\partial z_j}
 =\frac1\tau p_k(\delta_{kj}-p_j),\qquad
 \frac{\partial\log p_k}{\partial z_j}
 =\frac1\tau(\delta_{kj}-p_j),
 \label{eq:kd-softmax-derivative}
\end{equation}
其中 $\delta_{kj}$ 在 $k=j$ 时为一，否则为零。将其代入固定目标 $\mathbf q$ 的交叉熵，可得
\begin{equation}
 \frac{\partial H(\mathbf q,\mathbf p)}{\partial z_j}
 =-\sum_kq_k\frac{\delta_{kj}-p_j}{\tau}
 =\frac{p_j-q_j}{\tau}.
 \label{eq:kd-soft-gradient}
\end{equation}
第二步使用了 $\sum_kq_k=1$。这说明学生对某类分配的概率过高时，该类 logit 的梯度为正，梯度下降会降低它；概率过低时则相反。所有类别的梯度之和为零，与 softmax 的平移不变性一致。

温度为一、目标为硬标签时，梯度变成 $p_j-\mathbf1\{j=y\}$。正确类别被推高，其他类别被推低；但在两个错误类别之间，硬标签只通过学生自己的概率决定推动幅度，没有直接提供教师所学到的相似关系。软标签把这种关系放入 $q_j$ 中。不过，来自教师的关系也可能包含错误偏好，所以“多了一个监督信号”并不等于“多了一个正确监督信号”。

\subsection{网络容量与结构约束}
学生不是为每个输入单独存储一个任意概率向量，而是由共享参数生成所有输出。若它只能表示线性决策边界，就无法在所有输入上复现高度弯曲的教师边界。\keyterm{假设类（hypothesis class）} $\mathcal P_S=\{p_\theta:\theta\in\Theta\}$ 表示学生能够实现的函数集合。蒸馏求解的通常是教师在该集合中的某种近似投影；投影结果依赖数据分布、损失形式和优化过程，不只依赖学生参数个数。

参数较少的模型也不必在每个目标任务上弱于教师。如果任务只涉及教师能力的一小部分，小模型可能足以表示所需映射；教师还可能带有对当前任务无益的变化。另一方面，训练数据和目标不能凭空消除表示限制。讨论“学生超越教师”时，应说明比较的是哪项任务指标、哪个测试分布以及采用了什么额外监督，而不能把局部实验结果解释为全面继承且增强所有能力。

\section{经典输出蒸馏的目标与推导}
\label{sec:kd-classical}

\subsection{教师、学生和混合损失}
设固定教师分数为 $\mathbf a_\phi(x)$，学生分数为 $\mathbf z_\theta(x)$。在同一温度下定义
\[
 \mathbf q^\tau(x)=\operatorname{softmax}_\tau(\mathbf a_\phi(x)),\qquad
 \mathbf p^\tau(x)=\operatorname{softmax}_\tau(\mathbf z_\theta(x)).
\]
若训练样本都有真实标签，一种经典混合目标为
\begin{equation}
 J(\theta)=\frac1n\sum_{i=1}^n\left[
 (1-\alpha)H(\mathbf e_{y_i},\mathbf p^1(x_i))
 +\alpha\tau^2D_{\mathrm{KL}}(\mathbf q^\tau(x_i)\|\mathbf p^\tau(x_i))
 \right]+\lambda\Omega(\theta),
 \label{eq:kd-classical}
\end{equation}
其中 $\alpha\in[0,1]$ 为软监督权重，$\lambda\geq0$ 为正则项权重，$\Omega$ 为指定的参数正则函数。硬标签项通常在温度一计算，软标签项在教师和学生的相同温度下计算。若只具有无标签数据，可以保留软监督项，但没有真实标签就失去了直接纠正教师错误的一条信息来源。

式\eqref{eq:kd-classical}中的 $\tau^2$ 是常见梯度尺度补偿，原始温度蒸馏工作对此进行了讨论\cite{hinton2015}；其作用条件将在第\ref{sec:kd-temperature}节推导。它不是保证效果的定理，也不使不同温度的训练完全相同。不同论文可能把硬、软项分别乘以两个独立权重，也可能省略平方补偿后重新调节权重。比较方法前必须先把损失定义、平均方式和温度约定统一，否则相同数字的超参数没有可比性。

固定教师时，不考虑正则项，单样本混合损失对学生第 $j$ 个 logit 的梯度为
\begin{equation}
 g_j=(1-\alpha)(p_j^1-\mathbf1\{j=y\})
       +\alpha\tau(p_j^\tau-q_j^\tau).
 \label{eq:kd-total-gradient}
\end{equation}
教师无需接受反向传播，学生仍通过式\eqref{eq:kd-chain}更新自身参数。如果硬标签与教师明显冲突，两项可能在一些参数方向上相互抵消。这不是程序必然出错，而是混合目标确实要求在两个信号间折中；应进一步检查教师是否出错、样本是否有歧义，以及权重是否合适。

\subsection{温度为一时的混合标签解释}
\begin{proposition}[固定样本的有效目标]
设 $\tau=1$，教师固定，$\alpha\in[0,1]$，并允许学生输出在整个概率单纯形内自由选择。忽略与学生无关的常数后，混合损失等于以
\begin{equation}
 \mathbf r_{\mathrm{mix}}=(1-\alpha)\mathbf e_y+\alpha\mathbf q
 \label{eq:kd-mixed-target}
\end{equation}
为目标的交叉熵，其最小值在 $\mathbf p=\mathbf r_{\mathrm{mix}}$ 处取得。
\end{proposition}
\begin{proof}
由式\eqref{eq:kd-ce-kl}，去掉常数 $-\alpha H(\mathbf q)$ 后，两项对 $-\log p_k$ 的系数相加即为 $r_{\mathrm{mix},k}$。再用 KL 非负性可知交叉熵在两个分布相同时最小。
\end{proof}

若有效目标含零分量，而学生使用有限 logit 的 softmax，则边界分布只能作为极限逼近，不一定由有限参数精确实现；若学生同时拟合多个输入，参数共享也可能阻止逐点达到这个最小值。温度不为一时，硬、软项使用不同分布 $\mathbf p^1$ 和 $\mathbf p^\tau$，因而不能不加说明地合并成式\eqref{eq:kd-mixed-target}。

\begin{example}[教师犯错时标签权重的作用]
设真实标签为第一类，教师输出 $\mathbf q=(0.1,0.8,0.1)$，取 $\tau=1$。当 $\alpha=0.5$ 时有效目标为 $(0.55,0.40,0.05)$，最大分量仍对应真实类别；当 $\alpha=0.8$ 时有效目标为 $(0.28,0.64,0.08)$，最优自由输出转而预测第二类。这里的结论只针对该点的理想化优化，但足以说明更强蒸馏权重可能扩大教师错误，而不是必然改善学生。
\end{example}

\subsection{标签平滑与暗知识}
\keyterm{标签平滑（label smoothing）}常用 $(1-\varepsilon)\mathbf e_y+\varepsilon\mathbf u$ 代替硬标签，其中 $\mathbf u=(1/K,\ldots,1/K)$，$\varepsilon\in[0,1]$。它向所有类别加入规定的概率质量；教师软标签通常随输入变化，能表达更具体的类别关系。二者都能降低极端监督的程度，但只有在教师分布恰为相应固定分布等特殊条件下，目标才会一致。

\keyterm{暗知识（dark knowledge）}通常指教师在非最大概率类别上表达的相对偏好。它不是可直接观察的“真实知识含量”，而是对输出结构的形象称呼。若次要类别的概率主要来自噪声、错误相关性或训练不足，保留它们也会转移这些问题。检验暗知识是否有用，可以比较完整软标签、仅保留教师最大类别、标签平滑，以及在保持最大类别不变时打乱其余类别概率等对照，而不能仅凭软分布更复杂就断言其有效。

\section{温度缩放、高温极限与 logit 匹配}
\label{sec:kd-temperature}

\subsection{温度如何改变分布}
从式\eqref{eq:kd-softmax}可直接得到
\begin{equation}
 \frac{p_i^\tau}{p_j^\tau}=\exp\left(\frac{z_i-z_j}{\tau}\right).
 \label{eq:kd-odds}
\end{equation}
温度增加时，任意固定分数差对应的概率比趋近一；温度减小时，较大分数的优势被放大。有限分数下，$\tau\to\infty$ 时概率趋于均匀分布；$\tau\to0^+$ 时概率集中到最大分数对应的类别，若最大值并列，则极限在所有并列最大类别上均匀分配。正温度不会改变分数的次序，所以单独对一个固定模型做温度缩放不会改变其最大概率类别。

对固定 $\mathbf z$，记 $Z(\tau)=\sum_ke^{z_k/\tau}$、$\mu(\tau)=\sum_kp_k^\tau z_k$。由 $H(\mathbf p^\tau)=\log Z(\tau)-\mu(\tau)/\tau$ 和 $\mu'(\tau)=-\operatorname{Var}_{\mathbf p^\tau}(z)/\tau^2$，得到
\begin{equation}
 \frac{dH(\mathbf p^\tau)}{d\tau}
 =\frac{\operatorname{Var}_{\mathbf p^\tau}(z)}{\tau^3}\geq0.
 \label{eq:kd-entropy-temperature}
\end{equation}
方差指把类别索引按 $\mathbf p^\tau$ 抽样后，分数 $z_k$ 的方差。该等式由有限和求导得到，无需交换无限极限。它精确说明升温会软化固定模型的分布，却没有声称升温会使模型更准确或更校准。

\begin{example}[温度带来的软化]
取 $\mathbf z=(2,1,0)$。温度一时，概率约为 $(0.665241,0.244728,0.090031)$；温度二时，约为 $(0.506480,0.307196,0.186324)$。第二类相对第三类的概率比从 $e$ 变成 $e^{1/2}$，三个类别的顺序保持不变。温度增加使次要类别在损失中获得更可见的权重，但也削弱了主要类别的优势。
\end{example}

\subsection{为什么常乘以温度的平方}
给定一个输入，定义中心化分数 $\widetilde z_k=z_k-\bar z$、$\widetilde a_k=a_k-\bar a$，其中 $\bar z=K^{-1}\sum_kz_k$，$\bar a$ 同理。固定这些有限分数，让 $\tau\to\infty$，利用指数函数展开和中心化分量和为零，得到
\begin{equation}
 p_k^\tau=\frac1K+\frac{\widetilde z_k}{K\tau}+O(\tau^{-2}),\qquad
 q_k^\tau=\frac1K+\frac{\widetilde a_k}{K\tau}+O(\tau^{-2}).
 \label{eq:kd-high-temperature}
\end{equation}
于是未乘补偿的软损失梯度满足
\begin{equation}
 \frac{\partial D_{\mathrm{KL}}(\mathbf q^\tau\|\mathbf p^\tau)}{\partial z_k}
 =\frac{p_k^\tau-q_k^\tau}{\tau}
 =\frac{\widetilde z_k-\widetilde a_k}{K\tau^2}+O(\tau^{-3}).
 \label{eq:kd-tau-squared}
\end{equation}
一个 $1/\tau$ 来自 softmax 求导，另一个来自两分布差异在高温下缩小。乘以 $\tau^2$ 后，主导梯度规模不再按 $\tau^{-2}$ 衰减。这个结论要求温度相对于中心化分数足够大；有限温度下分布差异的变化不必精确服从此比例。因此，补偿后仍需根据验证集选择软项权重。

\subsection{高温 KL 与中心化平方误差}
还可以直接在均匀分布附近对 KL 展开。设 $u=1/K$，$q_k=u+\delta_k$、$p_k=u+\epsilon_k$，且两组扰动之和均为零。对对数作二阶展开并求和，一阶项抵消，得到
\[
 D_{\mathrm{KL}}(\mathbf q\|\mathbf p)
 =\frac1{2u}\sum_k(\delta_k-\epsilon_k)^2
 +O\!\left((\|\boldsymbol\delta\|_2+\|\boldsymbol\epsilon\|_2)^3\right).
\]
当 $K$ 固定且扰动足够小，使所有分量保持正数时，这一局部展开有效。代入式\eqref{eq:kd-high-temperature}，得到
\begin{equation}
 \lim_{\tau\to\infty}\tau^2D_{\mathrm{KL}}(\mathbf q^\tau\|\mathbf p^\tau)
 =\frac1{2K}\sum_{k=1}^K(\widetilde a_k-\widetilde z_k)^2.
 \label{eq:kd-logit-limit}
\end{equation}
这给出高温匹配相对分数的精确极限。需要中心化的原因是概率不识别共同平移：教师分数 $(4,2,0)$ 和学生分数 $(104,102,100)$ 给出相同概率，但直接对未中心化分数作平方匹配仍会产生很大误差。

高温下概率接近均匀，并不意味着乘以 $\tau^2$ 的目标完全失去信息；式\eqref{eq:kd-logit-limit}说明一阶概率差仍携带相对分数。不过温度很高时，原先概率极低的类别也会在中心化平方目标中受到较强约束，而这些分数可能缺乏可靠训练。温度太低则接近硬伪标签，细微关系较难参与学习。因此温度决定了强调何种差异，最优温度取决于教师、学生和任务，不能仅由“越软越好”确定。

\subsection{蒸馏温度与校准温度的区别}
蒸馏温度在训练期间改变监督几何；\keyterm{温度校准（temperature scaling）}则通常在独立验证集上为一个已训练模型选择缩放温度，使其负对数似然等概率指标改善。两者操作形式相似，但目的和使用阶段不同。学生部署时一般恢复温度一；如果还需校准，应另用验证数据估计。不能直接把教师训练时采用的蒸馏温度当成学生部署概率已经校准的证据。

\section{蒸馏为何可能有效，以及为何可能失败}
\label{sec:kd-mechanisms}

\subsection{条件期望、偏差与方差的简化分析}
固定输入 $x$，若教师恰好输出真实条件分布 $\mathbf r(x)$，那么独热随机标签的条件期望满足 $\mathbb E[\mathbf e_Y\mid X=x]=\mathbf r(x)$。在温度一、参数固定时，硬标签 logit 梯度为 $\mathbf p(x)-\mathbf e_Y$，其条件期望为 $\mathbf p(x)-\mathbf r(x)$，恰好等于使用理想软标签的梯度。软标签相当于直接给出这个条件平均，从而消除了由抽取单个标签导致的条件随机波动。

更具体地，固定教师输出 $\mathbf q(x)$、设混合目标 $\mathbf U=(1-\alpha)\mathbf e_Y+\alpha\mathbf q(x)$，则
\begin{align}
 \mathbb E[\mathbf U\mid X=x]-\mathbf r(x)
 &=\alpha(\mathbf q(x)-\mathbf r(x)),\label{eq:kd-target-bias}\\
 \operatorname{Cov}(\mathbf U\mid X=x)
 &=(1-\alpha)^2\left[\operatorname{diag}(\mathbf r(x))-
 \mathbf r(x)\mathbf r(x)^\top\right].\label{eq:kd-target-variance}
\end{align}
这里 $\operatorname{diag}(\mathbf r)$ 表示以 $\mathbf r$ 为对角的矩阵。第一式揭示教师误差引入的条件偏差，第二式揭示标签抽样噪声的减少。它们是在教师固定、只考虑新标签随机性的简化模型中成立的恒等式；不包括教师由有限训练集学习产生的随机性，也不等于已经证明学生测试误差降低。

\subsection{归纳偏置与函数空间约束}
教师能够把许多训练样本中学到的结构压缩成输入相关的目标，使学生获得比单个标签更细的监督。另一方面，匹配教师也在学生函数空间中施加额外约束，这与\keyterm{正则化（regularization）}的思想相通：减少某些不受约束的变化，可能降低过拟合，也可能排除真实任务需要的变化。蒸馏不是简单增加样本标签数量，而是改变哪些函数被训练目标偏好。

如果学生容量足够、蒸馏输入覆盖部署分布、教师可靠且优化充分，模仿教师可能带来良好结果；这些条件任一不满足都会影响迁移。教师只在常见样本上正确，学生可能继承长尾错误；教师使用训练中的背景捷径，学生也可能模仿这种相关性；学生过小，可能只能保留易拟合但无用的教师行为。因此教师选择不应仅看一个总体准确率，还应看错误类型、置信度、任务覆盖和与学生架构的兼容性。

\subsection{一个有限而清晰的误差界}
\begin{proposition}[对教师的一致性只能控制一部分错误]
设教师和学生在同一输入空间上给出确定类别预测 $h_T,h_S$，真实标签为 $Y$。在任意共同评估分布上，有
\begin{equation}
 P(h_S(X)\ne Y)\leq P(h_T(X)\ne Y)+P(h_S(X)\ne h_T(X)).
 \label{eq:kd-disagreement-bound}
\end{equation}
\end{proposition}
\begin{proof}
若学生出错而教师正确，则二者预测必定不同。因此学生出错事件包含在“教师出错”与“师生不同”两事件的并集中，对并集使用概率的次可加性即可。
\end{proof}

这个界直接展示两项来源：教师本身的错误和学生对教师的模仿误差。它不保证学生一定不如教师，也不说明训练集上的分布匹配已经控制了部署时的类别分歧。特别是教师在两个类别之间几乎打平时，极小的概率变化就可能改变最大类别。要从概率误差进一步控制类别分歧，需要教师决策间隔等附加条件。

例如教师最大概率类别为 $k^*$，与第二大概率之差为 $\gamma>0$。若在该输入上 $\|\mathbf p-\mathbf q\|_\infty<\gamma/2$，则对任意 $j\ne k^*$，$p_{k^*}-p_j>\gamma-2(\gamma/2)=0$，学生必与教师同类。间隔很小时要求更精细的概率逼近。这也说明，仅报告平均 KL 而不观察决策边界附近样本，可能遗漏重要的分类错误。

\section{响应蒸馏与目标类别的解耦}
\label{sec:kd-response}

\subsection{输出中究竟有哪些信息}
\keyterm{响应蒸馏（response-based distillation）}直接比较任务输出，包括分类概率、logit、回归值或结构化预测。在分类中，完整分布同时包含“正确类别相对所有其余类别的优势”和“其余类别之间的相对关系”。经典 KL 把二者放在同一个目标中，因而它们的相对权重不是完全自由的。分析这种耦合，有助于理解某些教师非常自信时，非目标类别关系为何难以影响训练。

固定一个真实类别 $y$，假设 $q_y,p_y\in(0,1)$。定义二元分布 $\mathbf b_q=(q_y,1-q_y)$、$\mathbf b_p=(p_y,1-p_y)$，以及排除目标类别后的条件分布 $\bar q_k=q_k/(1-q_y)$、$\bar p_k=p_k/(1-p_y)$，其中 $k\ne y$。将每个非目标概率写成“总质量乘以条件概率”，可得
\begin{equation}
 D_{\mathrm{KL}}(\mathbf q\|\mathbf p)
 =D_{\mathrm{KL}}(\mathbf b_q\|\mathbf b_p)
 +(1-q_y)D_{\mathrm{KL}}(\bar{\mathbf q}\|\bar{\mathbf p}).
 \label{eq:kd-decomposition}
\end{equation}
推导只需把非目标部分中的对数拆成 $\log[(1-q_y)/(1-p_y)]+\log(\bar q_k/\bar p_k)$，再利用条件概率之和为一。第一项关注目标类别与其余类别的二分关系，第二项关注其余类别内部的相对分配。

\subsection{解耦的作用与边界}
\keyterm{解耦知识蒸馏（decoupled knowledge distillation, DKD）}利用这种分解，为目标类别项和非目标类别项配置独立权重\cite{zhao2022}。其思想可概括为
\[
 L_{\mathrm{dec}}=\beta_1D_{\mathrm{KL}}(\mathbf b_q\|\mathbf b_p)
 +\beta_2D_{\mathrm{KL}}(\bar{\mathbf q}\|\bar{\mathbf p}),\qquad \beta_1,\beta_2\geq0.
\]
它是对原目标的重新设计，而不是在任意权重下都与经典 KL 相等。若教师对目标类别给出 $0.999$ 的概率，原目标中非目标关系只乘以 $0.001$；独立加权可以使该部分更显著，但也可能放大教师不可靠的尾部预测。

教师与学生应在同一温度下形成用于分解的分布，温度补偿再按具体目标约定处理。若类别数只有二，排除真实类别后只剩一个类别，条件分布恒为 $(1)$，非目标 KL 恒为零，解耦便不会凭空产生额外关系。若没有真实标签而使用教师最大类别代替 $y$，就构成另一种算法假设，应单独说明，不能把它当成已知真实类别的原设置。

\section{特征蒸馏：对齐中间表示}
\label{sec:kd-feature}

\subsection{为什么需要中间层监督}
最终输出把输入压缩成与任务有关的少数数值。教师中间层可能还保留边缘、局部组合、语义结构或词元上下文等信息，学生可以通过匹配这些表示获得更多训练约束。\keyterm{特征蒸馏（feature-based distillation）}因此把监督放到网络内部。FitNets 使用教师中间表示为较薄的学生提供提示，是这一方向的代表工作\cite{romero2015}。中间监督可能改善优化，也可能限制学生使用更适合自身架构的表示；它并不比输出蒸馏无条件更强。

设教师第 $l$ 层特征为 $\mathbf F_T^{(l)}$，学生被选中的第 $m(l)$ 层特征为 $\mathbf F_S^{(m(l))}$。如果维度不同，可引入可训练适配器 $A_l$，构造
\begin{equation}
 L_{\mathrm{feat}}=\sum_{l\in\mathcal I}\frac{\omega_l}{N_l}
 \left\|A_l(\mathbf F_S^{(m(l))})-
       \operatorname{sg}(\mathbf F_T^{(l)})\right\|_F^2.
 \label{eq:kd-feature-loss}
\end{equation}
其中 $\mathcal I$ 为所选教师层的集合，$\omega_l\geq0$ 为层权重，$N_l$ 为该层参与比较的标量数。适配器可改变通道维度；空间尺度不同还需明确插值、池化或其他位置对应操作。除以 $N_l$ 避免仅因某层元素更多就使其损失数值更大，但不同层的激活尺度仍可能不同。

\subsection{对齐不只是把维度改成一样}
两个张量形状相同，不意味着其通道语义相同。将隐藏通道重排，并对后续层做相应逆重排，网络整体函数可以保持不变，但逐通道平方误差会发生变化。更一般地，某些可逆基变换也能被后一线性层抵消。因此直接特征匹配既约束任务信息，也约束教师选用的表示坐标。适配器减轻这种坐标差异，但若适配器过强，也可能由它承担大部分拟合，而学生主体没有获得期望能力。

选择对应层时，应考虑感受野、空间分辨率和语义抽象层次，而不只按层编号等比例连接。\keyterm{感受野（receptive field）}指一个特征位置在输入中能够依赖的区域。学生浅层局部特征与教师深层全局特征即使维度一致，其信息可达范围也可能不同。实际设计可先选择少量合理对应层，建立有效基线，再评估更多层监督是否有收益；无区别地匹配所有层容易引入冗余约束和额外开销。

若适配器只用于辅助训练分支，部署时可以移除；若它位于学生主预测路径中，就必须保留并计入推理成本。报告模型大小时不能只计算学生主干而忽略推理必需的适配模块。特征蒸馏通常还要求在教师前向时保存所选激活，即使没有教师反向传播，峰值内存也可能显著高于仅匹配最终输出的方案。

\subsection{注意力迁移与归一化}
卷积特征可通过通道聚合形成空间图，例如对单样本 $\mathbf F\in\mathbb R^{C\times H\times W}$ 定义 $A(\mathbf F)_{hw}=\sum_c|F_{chw}|^2$。把此图展开为向量，再用 $\widehat{\mathbf a}=\mathbf a/\max(\|\mathbf a\|_2,\epsilon)$ 归一化，其中 $\epsilon>0$ 用于零范数保护，便可比较师生在哪些区域有较强响应。注意力迁移的原始研究讨论了此类监督\cite{zagoruyko2017}。这里的“注意力”是由激活汇总构造的空间强度，并不等同于 Transformer 的注意力权重，也不是严格的因果解释。

归一化使匹配更关注相对空间分布，弱化整体尺度差异，但会丢失幅值信息。若所有激活接近零，归一化和保护常数的影响可能主导目标。应检查教师、学生特征范数分布，而不只观察归一化后的图是否相似。对于裁剪或翻转后的输入，空间位置还必须对应；在不同视图之间逐像素匹配而不处理几何变换，会把正确的局部特征当成错误。

\section{关系蒸馏与对比学习基础}
\label{sec:kd-relation}

\subsection{从单个表示转向样本之间的关系}
\keyterm{关系蒸馏（relational knowledge distillation）}比较多个样本在表示空间中的相互关系。设一批样本的教师向量为 $\mathbf h_i^T\in\mathbb R^{d_T}$，学生向量为 $\mathbf h_i^S\in\mathbb R^{d_S}$。即使 $d_T\ne d_S$，仍可比较距离 $\|\mathbf h_i^T-\mathbf h_j^T\|_2$ 与 $\|\mathbf h_i^S-\mathbf h_j^S\|_2$。Park 等人的工作研究了距离与角度关系的迁移\cite{park2019}，说明可传递的对象不必是单个向量坐标。

为弱化整体尺度，记批内非对角距离均值为 $\mu_T=\frac{2}{B(B-1)}\sum_{i<j}\|\mathbf h_i^T-\mathbf h_j^T\|_2$，学生的 $\mu_S$ 同理。若两者均正，可定义归一化距离 $d_{ij}^T=\|\mathbf h_i^T-\mathbf h_j^T\|_2/\mu_T$，并优化
\begin{equation}
 L_{\mathrm{rel}}=\frac{2}{B(B-1)}\sum_{i<j}
 \rho(d_{ij}^S-d_{ij}^T),
 \label{eq:kd-relation-loss}
\end{equation}
其中 $\rho$ 可以是平方函数或具有较弱异常值敏感性的稳健损失。若所有特征相同，距离均值为零，应规定保护规则；这时表示已经发生塌缩，不能仅通过除数加一个小常数就认为问题被解决。

欧氏距离对共同平移和正交变换保持不变，归一化后还对正比例缩放保持不变。这让关系匹配比逐坐标匹配更自由，但也意味着只凭这些关系不能唯一恢复教师坐标系。进一步比较三点之间的夹角，可以约束局部几何结构；其计算涉及更多元组，成本也更高。关系目标通常和任务目标结合，避免学生仅复现某种几何排列，却没有学到正确类别决策。

\subsection{Gram 矩阵与相似性保持}
把归一化后的表示按行组成 $\mathbf H_T\in\mathbb R^{B\times d_T}$、$\mathbf H_S\in\mathbb R^{B\times d_S}$，Gram 矩阵 $\mathbf G_T=\mathbf H_T\mathbf H_T^\top$ 和 $\mathbf G_S=\mathbf H_S\mathbf H_S^\top$ 都属于 $\mathbb R^{B\times B}$。若每行具有单位范数，矩阵元素就是样本间余弦相似度。匹配 $\|\mathbf G_S-\mathbf G_T\|_F^2$ 可以绕开隐藏维度不一致的问题，但其存储规模为 $O(B^2)$，批量很大时并不便宜。

批次组成决定哪些关系被看见。若一批全部来自同一类别，便难以提供跨类几何信息；若随机抽到大量近重复样本，关系目标也可能重复强调少数方向。扩大批量、维护历史表示或采样特定样本对都能改变关系覆盖，但会引入计算成本、表示陈旧性或采样偏差。不能把采用了关系目标本身当成已经获得全数据几何结构的证据。

\subsection{对比蒸馏的正负样本}
\keyterm{对比学习（contrastive learning）}通过比较正样本对与负样本对学习表示。一个教学用批内目标为：同一输入的学生表示 $\mathbf u_i$ 与教师表示 $\mathbf v_i$ 构成正对，不同输入构成候选负对，所有表示经过适配并归一化到相同维度，定义
\begin{equation}
 L_{\mathrm{ctr}}=-\frac1B\sum_{i=1}^B
 \log\frac{\exp(\mathbf u_i^\top\mathbf v_i/\kappa)}
 {\sum_{j=1}^B\exp(\mathbf u_i^\top\mathbf v_j/\kappa)},\qquad \kappa>0.
 \label{eq:kd-contrastive}
\end{equation}
这里 $\kappa$ 是对比温度，不必等于分类蒸馏温度 $\tau$。损失要求学生从一批教师表示中辨认自己的对应项。对比表示蒸馏的代表研究为 CRD\cite{tian2020}；式\eqref{eq:kd-contrastive}用于解释基本机制，不是对该论文全部采样方式和损失实现的逐项复现。

不同样本可能具有同一语义，因而“不同输入必为负例”是一种假设。错误负例会迫使本应接近的表示分离。某些对比目标在特定联合分布、负样本独立采样等假设下与互信息下界有关，但不能由任意对比损失值直接读出真实互信息，更不能据此证明迁移了全部任务知识。实际评价仍需看下游预测和分布变化下的行为。

\section{多教师、在线蒸馏与自蒸馏}
\label{sec:kd-training-schemes}

\subsection{多教师的组合与分歧}
设有 $M$ 个教师，输出分布分别为 $\mathbf q^{(m)}(x)$。给定非负权重 $w_m$ 且 $\sum_mw_m=1$，算术混合 $\bar{\mathbf q}=\sum_mw_m\mathbf q^{(m)}$ 是合法概率分布。固定教师时，$\sum_mw_mH(\mathbf q^{(m)},\mathbf p)=H(\bar{\mathbf q},\mathbf p)$，因此加权交叉熵可以解释为匹配平均分布。KL 版本对学生的梯度同样一致，但数值会相差与学生无关的教师熵项。

平均 logit 再做 softmax 与平均概率一般不同。若各教师使用同一温度且输出正概率，对 logit 加权平均得到的分布正比于 $\prod_m(q_k^{(m)})^{w_m}$，属于归一化几何组合。几何组合对某教师接近零的概率更敏感，算术组合则保留各教师分配的质量。没有一种组合规则能在不了解任务与教师误差的情况下总是最好。若各教师的类别顺序或标签定义不同，必须先建立共同语义空间，否则求平均没有正确含义。

教师分歧既可能提供互补信息，也可能暴露不可靠监督。依照输入选择教师的门控权重可以更灵活，但门控本身需要可靠训练和验证。若权重依赖学生且参与求导，便不能继续把组合教师视为固定目标；若只希望它决定样本权重，应明确停止相应梯度。多教师带来的准确率变化还必须与额外推理、存储和生成目标成本一起报告。

\subsection{离线、在线和相互学习}
\keyterm{离线蒸馏（offline distillation）}先训练教师，再固定教师训练学生；目标可以现场计算，也可以缓存。“离线”描述教师训练与学生训练的组织关系，并不意味着一定把预测预先存入硬盘。\keyterm{在线蒸馏（online distillation）}则在训练过程中产生或更新教师信号，例如多个网络共同学习。深度相互学习\cite{zhang2018}让模型在标签监督之外彼此匹配输出，提供了不预先训练单个强教师的方案。

两个网络可分别优化 $L_1=H(\mathbf e_y,\mathbf p_1)+\beta D_{\mathrm{KL}}(\operatorname{sg}(\mathbf p_2)\|\mathbf p_1)$ 和相应的 $L_2$。更新网络一时把网络二的目标视为常量，更新网络二时反过来处理。这不同于把两个分布都接入同一个可微 KL 然后同时求导。共同训练的模型可能互相补充，也可能迅速强化相同错误；标签、数据增强、初始化差异等会影响它们是否保留有益多样性。

\subsection{同结构再训练与指数移动平均教师}
\keyterm{自蒸馏（self-distillation）}包含多种组织方式：用一个已训练模型指导同结构的新模型，用深层分支指导浅层分支，或用历史模型指导当前模型。Born Again Networks 研究了与教师同结构的学生再训练\cite{furlanello2018}。此时即使学生取得更好的某项性能，额外训练阶段与优化差异也必须计入解释，不能把效果全部归因于结构压缩。

另一种常见教师由学生参数的\keyterm{指数移动平均（exponential moving average, EMA）}形成：
\begin{equation}
 \phi_t=\mu\phi_{t-1}+(1-\mu)\theta_t,\qquad 0\leq\mu<1.
 \label{eq:kd-ema}
\end{equation}
Mean Teacher 将此类教师用于一致性训练\cite{tarvainen2017}。递推展开为 $\phi_t=\mu^t\phi_0+(1-\mu)\sum_{j=1}^t\mu^{t-j}\theta_j$，说明近期参数权重较大。EMA 平滑参数轨迹，但神经网络是非线性的，平均参数的预测不等于历史模型预测的平均；也不能由参数更平滑就断定教师总比学生更准确。

\subsection{无标签自监督与塌缩问题}
\keyterm{自监督学习（self-supervised learning）}从数据自身构造目标。若师生分别观察同一图像的不同裁剪，要求输出相近，可以鼓励对这些变换稳定的表示。然而只要求两者一致存在平凡解：所有输入都输出同一个向量。这个现象称为\keyterm{表示塌缩（representation collapse）}。停止教师梯度可以改变训练动力学，但单独停止梯度并不是排除所有塌缩解的数学保证。

DINO 将自蒸馏用于无标签视觉表示学习\cite{caron2021}，其设计包含移动平均教师、不同视图及对输出分布的处理。理解这一方向时，重点应放在目标如何同时鼓励跨视图一致与保留跨样本区分，而不是把任何无标签 KL 匹配都称为同一种算法。验证表示质量还需要指定线性探测、近邻检索或下游微调协议；仅看蒸馏损失下降不能排除学生已经学会一个无用常量。

\section{蒸馏数据、分布偏移与无原始数据蒸馏}
\label{sec:kd-data}

\subsection{在什么输入上提问，决定学到什么}
固定教师后的总体蒸馏目标可以写成
\begin{equation}
 J_Q(\theta)=\mathbb E_{X\sim Q_X}
 D_{\mathrm{KL}}(\mathbf q(X)\|\mathbf p_\theta(X)).
 \label{eq:kd-input-distribution}
\end{equation}
如果 $Q_X$ 只覆盖白天道路，目标便不直接约束夜间道路上的师生差异。学生可能依靠结构归纳偏置迁移到夜间，但这需要验证。一个模型在蒸馏数据上的损失接近零，只说明在该分布加权的区域接近教师，不意味着在全部可能输入上复现教师。对于开放式任务，输入采样策略与损失形式同样重要。

若部署输入分布 $P_X^{\mathrm{deploy}}$ 相对于 $Q_X$ 绝对连续，并且知道密度比 $w(x)=dP_X^{\mathrm{deploy}}/dQ_X$，则对可积损失有 $\mathbb E_{P_X^{\mathrm{deploy}}}\ell(X)=\mathbb E_{Q_X}[w(X)\ell(X)]$。这是\keyterm{重要性加权（importance weighting）}的基础。它要求部署分布不存在蒸馏分布完全不支持的区域；即使条件满足，密度比很大也会使估计方差很高。实际通常不知道精确密度比，因此补充覆盖数据和单独评估分布变化，往往比形式上写出权重更关键。

\subsection{数据增强、一致性与可靠性选择}
\keyterm{数据增强（data augmentation）}通过变换 $v(x)$ 产生训练视图。若变换保持标签语义，可让师生看到同一视图；也可以让教师看到较弱增强、学生看到较强增强，训练二者输出一致。后者还引入了“该变换不改变任务答案”的假设。将关键物体完全裁掉、把数字翻转成另一数字，或删去决定文本情感的否定词，都可能违反这个假设。增强不正确时，更强一致性会强化错误监督。

教师置信度可以用于过滤样本或调节权重，但必须警惕选择偏差。只保留最高置信度样本可能主要留下容易样本，使边界、长尾类别和异常输入缺乏监督；教师对某些错误恰好非常自信时，过滤还会保留最危险的错误。一个较稳妥的实验应同时记录保留比例、类别分布、置信度与正确率的关系，并比较不过滤的基线。改变采样后，任何性能变化都可能来自数据分布改变，而非蒸馏损失本身。

缓存教师预测时，还要明确缓存的是原图、某次随机增强还是多个固定视图。学生每次看到随机新裁剪，却总匹配原图的局部特征，可能造成目标错位。标签保真的全局输出一致性与逐位置特征对应是不同约束，不能用同一个“增强不变”说法笼统处理。缓存文件应记录教师版本、类别顺序、预处理、温度和样本标识，避免把旧目标错误地接到新输入上。

\subsection{无原始数据蒸馏的可用信息}
\keyterm{无原始数据蒸馏（data-free knowledge distillation）}指无法访问教师的原始训练数据，并不意味着完全没有任何信息或先验。可用资源可能包括教师权重、模型中的归一化统计量、额外公开输入、生成器或少量任务描述。不同设置的信息条件差异很大，应先说明可以访问什么，再比较方法效果。使用大量额外公开数据的方案与只访问教师参数的方案，不是同等信息预算。

以可微教师为例，可以优化合成输入 $x$，使其被教师判为指定类别，并对输入施加自然性或统计匹配约束。一个概念性目标是
\[
 L_{\mathrm{syn}}(x)=H(\mathbf e_c,\mathbf q(x))
 +\lambda_{n}\Omega_{n}(x)+\lambda_{s}\Omega_{s}(x),
\]
其中 $c$ 是指定类别，$\Omega_{n}$ 是输入先验，$\Omega_{s}$ 可度量激活统计与教师保存统计之间的差别。DeepInversion 使用模型中可用的统计信息进行数据反演，是这一方向的代表工作\cite{yin2020}。此处只是解释目标组成，具体层、正则项和优化细节应以原论文为准。

教师对合成输入高置信并不保证该输入位于真实数据分布中。优化器可能找到对教师有效、对真实任务无意义的模式。因此需要同时检查类别覆盖、样本多样性和真实测试集性能。对抗式合成还可以寻找师生差异最大的输入，但若没有合理输入约束，这种差异可能集中在无关区域。无原始数据蒸馏并不会消除数据分布问题，而是把它转化为如何构造有用查询的问题。

\section{语言模型的前置知识与序列蒸馏}
\label{sec:kd-language}

\subsection{词元、自回归分解与教师强制}
\keyterm{词元（token）}是分词器处理文本后产生的离散单位，可能是一段字节、一个字、一个词或子词，不等于固定数量的汉字。设词表为 $\mathcal V$，其大小为 $V$；输入提示为 $x$，输出词元序列为 $y=(y_1,\ldots,y_L)$。\keyterm{自回归模型（autoregressive model）}按条件概率链式法则写成
\begin{equation}
 p_\theta(y\mid x)=\prod_{t=1}^L p_\theta(y_t\mid x,y_{<t}),
 \qquad y_{<t}=(y_1,\ldots,y_{t-1}).
 \label{eq:kd-autoregressive}
\end{equation}
其中序列终止可以由结束词元表示。这里的 $t$ 是序列位置，不是前文的训练迭代索引；每次使用时由上下文区分。位置条件分布由网络计算，生成算法再决定如何从该分布选择下一个词元。

有参考序列时，\keyterm{教师强制（teacher forcing）}指训练每个位置的预测时，输入参考序列已经出现的正确前缀。这个名称中的“教师”不一定是蒸馏教师，也可能只是人工参考文本。推理时模型却看到自己生成的前缀，早先的错误会改变后续输入；这种训练与推理前缀分布的差别常称为\keyterm{暴露偏差（exposure bias）}。即使每个训练前缀上的损失较小，也可能在学生自行产生的陌生前缀上表现很差。

\subsection{词元级概率蒸馏}
教师和学生使用相同词表、同一上下文格式、同一前缀时，可在每个位置比较完整词表分布：
\begin{equation}
 L_{\mathrm{token}}=\frac{\tau^2}{\sum_{i,t}m_{it}}
 \sum_{i,t}m_{it}D_{\mathrm{KL}}\left(
 q^\tau(\cdot\mid x_i,y_{i,<t})\|
 p_\theta^\tau(\cdot\mid x_i,y_{i,<t})\right).
 \label{eq:kd-token-loss}
\end{equation}
其中 $m_{it}\in\{0,1\}$ 表示该位置是否参与损失，总有效位置数必须为正。提示部分、填充位置和不希望监督的位置应按目标明确屏蔽。用词元总数平均使每个词元等权；先对每条序列平均再对样本平均使每条序列等权，两者在序列长度不同时代表不同目标。

词元级 KL 传递的不只是最终写出的词元，还包括教师在当前前缀下对其他候选的偏好。但词表往往很大，保存每个位置的完整概率会产生 $O(BLV)$ 的目标存储量。只保留概率最大的若干候选可以节省资源，却改变了可观测信息；若教师接口仅返回生成文本，则无法精确恢复完整词表 KL。不能把没有获得的概率分布假定为已经蒸馏。

\subsection{序列级监督与生成数据}
\keyterm{序列级蒸馏（sequence-level distillation）}可以让教师先生成目标序列 $\widehat y(x)$，再以
\begin{equation}
 L_{\mathrm{seq}}(\theta)=-\sum_{t=1}^{|\widehat y|}
 \log p_\theta(\widehat y_t\mid x,\widehat y_{<t})
 \label{eq:kd-sequence-loss}
\end{equation}
训练学生。Kim 与 Rush 的工作系统研究了序列级知识蒸馏\cite{kim2016}。这一方式只需要教师文本，适用于不能访问内部概率的情况；但单条生成结果只展示一种轨迹，并没有提供每个位置所有候选的相对概率。

如果 $\widehat Y$ 确实按完整教师序列分布 $q(\cdot\mid x)$ 随机采样，那么 $-\log p_\theta(\widehat Y\mid x)$ 是序列交叉熵的无偏蒙特卡洛估计。在允许交换期望与求导的条件下，其梯度也是相应目标梯度的无偏估计；有限词表和固定最大长度、参数邻域内导数有界是足够的简单情形。若用贪心、束搜索、截断采样或筛选后的输出，样本来自另一分布，就不再是原教师完整分布交叉熵的无偏估计。

\keyterm{束搜索（beam search）}保留若干当前评分较高的前缀继续扩展，它近似寻找高分序列；\keyterm{截断采样（truncated sampling）}则在筛选候选后重新归一化并抽样。两者都改变数据中的多样性和模式覆盖。教师生成的文本往往还经过正确性筛选、格式约束或去重，所得目标更接近“教师加生成及筛选流程”的行为。记录这个流程与记录教师名称同样重要。

\subsection{分词器不同与推理过程蒸馏}
若师生分词器不同，即使两边词表大小碰巧相同，索引对应的事件也可能不同，不能直接对齐同一索引计算 KL。一段教师词元在学生侧可能拆成多个词元，条件前缀也不同。最直接的行为迁移是解码教师文本后，用学生分词器重新编码再做序列监督；更精细的概率对齐需要共同事件空间或专门对齐方法。仅把向量长度补齐不会解决语义不一致。

\keyterm{推理过程蒸馏（reasoning distillation）}使用教师给出的中间解释、步骤或轨迹训练学生。中间文本可以把复杂任务分解成较容易预测的局部目标，但可读解释不必忠实描述教师内部计算。最终答案正确也不保证每一步推导正确；相反，合法的另一条推导可能与教师措辞不同。数学、程序或结构化任务应尽量采用可执行检查、逐步约束或独立验证，并区分答案正确率、步骤有效性和文本风格一致性。

长推理文本还会改变训练 token 预算和部署成本。若学生学会生成更长解释，单次参数前向虽更便宜，完整回答却不一定更快。比较推理蒸馏时应匹配训练词元数、推理长度上限和采样次数；仅比较模型参数量不能说明整个问题求解过程已经加速。对不需要显示过程的任务，还可分别研究带过程监督与只监督答案的效果，而不预设越长的过程越有益。

\section{正向 KL、反向 KL 与学生分布上的蒸馏}
\label{sec:kd-onpolicy}

\subsection{两个方向分别惩罚什么}
正向 KL 为 $D_{\mathrm{KL}}(q\|p)=\mathbb E_q[\log q-\log p]$，教师高概率区域在目标中权重大；学生若在这些区域给出极小概率，会受到强惩罚。反向 KL 为 $D_{\mathrm{KL}}(p\|q)=\mathbb E_p[\log p-\log q]$，学生高概率区域权重大；学生若把质量放到教师极不认可的区域，会受到强惩罚。两者都在分布完全一致时为零，方向差异主要在学生受到容量限制、支持限制或优化限制时体现。

“正向 KL 覆盖模式，反向 KL 选择模式”是一种有条件的直觉。若学生可以精确表示教师，两者共同最优解就是教师，并不存在必须丢掉某个模式的结论。若教师有多个分离峰，而学生只能表示一个窄峰，正向目标会强烈反对漏掉教师主要区域，反向目标则可能倾向把学生质量集中到其中一个高概率区域。最终结果仍取决于模型族；把这个直觉说成任意任务上的严格定理是不准确的。

固定正概率教师 $\mathbf q$，学生 $\mathbf p=\operatorname{softmax}_1(\mathbf z)$ 时，对反向 KL 求导可得
\begin{equation}
 \frac{\partial D_{\mathrm{KL}}(\mathbf p\|\mathbf q)}{\partial z_j}
 =p_j\left(\log\frac{p_j}{q_j}
             -D_{\mathrm{KL}}(\mathbf p\|\mathbf q)\right).
 \label{eq:kd-reverse-gradient}
\end{equation}
推导时先对 $p_k$ 求导得到 $\log(p_k/q_k)+1$，再乘 softmax 的 Jacobian 并求和，常数一相互抵消。它明显不同于正向 KL 的 $p_j-q_j$。因此交换函数参数不是无关紧要的代码变动，而是改变了梯度和优化偏好。

\subsection{序列 KL 的链式分解}
为避免终止问题，先考虑有限最大长度 $L$，将结束后的状态补为吸收状态。设两个模型对相关序列赋正概率。将式\eqref{eq:kd-autoregressive}代入序列 KL 并交换有限求和，可得
\begin{equation}
 D_{\mathrm{KL}}(q(Y\mid x)\|p_\theta(Y\mid x))
 =\sum_{t=1}^L\mathbb E_{Y_{<t}\sim q}
 D_{\mathrm{KL}}\left(q(\cdot\mid x,Y_{<t})\|
 p_\theta(\cdot\mid x,Y_{<t})\right).
 \label{eq:kd-sequence-chain}
\end{equation}
证明来自 $\log(q(Y)/p(Y))=\sum_t\log(q(Y_t\mid Y_{<t})/p(Y_t\mid Y_{<t}))$，再对每个前缀条件化。关键在于外层前缀由教师分布产生。反向序列 KL 则对应学生前缀期望下的反向条件 KL。外层采样分布和内层散度方向是一同决定目标的两个组成部分。

因此，在人工参考前缀上做正向词元 KL、在教师生成前缀上做正向词元 KL、在学生生成前缀上做正向词元 KL，是三个不同目标。第一个由参考数据覆盖决定，第二个在相应条件下对应正向序列 KL，第三个针对学生经常访问的状态。不能因为内层公式相同，就把三者称为同一个序列散度的精确估计。

\subsection{在策略蒸馏与梯度边界}
\keyterm{在策略蒸馏（on-policy distillation）}让学生生成序列或前缀，再请求教师在这些状态上提供监督。它直接面向学生实际容易走到的状态，从而缓解只在理想参考轨迹上学习造成的偏差。Agarwal 等人的 GKD 研究了这类语言模型蒸馏\cite{agarwal2024}。教师在学生错误前缀上的输出仍是条件续写分布，不一定构成事实纠错指令；若前缀已经包含错误假设，还需要任务设计使教师反馈具有想要的含义。

用 $d_\theta(h\mid x)$ 表示学生产生前缀 $h$ 的分布，一个概念性目标为 $J(\theta)=\mathbb E_{h\sim d_\theta}D(q(\cdot\mid h)\|p_\theta(\cdot\mid h))$。设 $D_\theta(h)$ 为内部损失，在可微且允许交换求导与期望时，完整梯度为
\begin{equation}
 \nabla_\theta J
 =\mathbb E_{h\sim d_\theta}\left[
 \nabla_\theta D_\theta(h)
 +D_\theta(h)\nabla_\theta\log d_\theta(h)\right].
 \label{eq:kd-onpolicy-gradient}
\end{equation}
第二项来自采样分布也依赖参数。把生成前缀当成固定数据，只对词元损失反向传播，会省略这一项，因此不能无条件称为上述完整目标的无偏梯度。

一种清楚的算法表述是：第 $s$ 轮固定旧学生 $\theta_s$ 生成数据，接着优化代理目标 $J_s(\theta)=\mathbb E_{h\sim d_{\theta_s}}D_\theta(h)$，然后重新采样。对此固定采样分布的代理目标，条件损失梯度是恰当的；它与对当前分布依赖求全导数的问题不同。是否使用策略梯度、怎样控制方差、多久刷新数据，都属于具体算法设计，不能由“在策略”三个字自动确定。

\subsection{与强化学习的联系和区别}
\keyterm{强化学习（reinforcement learning）}通过环境奖励优化行动策略，语言生成也可被看成逐词元行动过程。蒸馏提供教师行为或分布监督，强化学习提供奖励目标，两者可以结合，但教师模仿并不必然等于奖励最大化。若教师并非最优策略，过强模仿会限制学生改进；若奖励信号稀疏，教师轨迹又可能提供有用初始化。比较二者时应明确教师反馈是概率、示范、偏好还是标量奖励，并分别说明优化目标。

\section{从分类扩展到结构预测、多模态与生成模型}
\label{sec:kd-applications}

\subsection{回归与多标签任务}
回归任务预测连续值，教师给出 $f_T(x)$，学生给出 $f_S(x)$，最直接的蒸馏项是 $(f_S(x)-f_T(x))^2$。若教师还能估计不确定性，可把其输出解释为概率分布，但这需要明确分布假设。例如教师、学生分别给出一维高斯分布 $\mathcal N(\mu_T,\sigma_T^2)$、$\mathcal N(\mu_S,\sigma_S^2)$，且方差严格为正，则
\begin{equation}
 D_{\mathrm{KL}}(T\|S)=\log\frac{\sigma_S}{\sigma_T}
 +\frac{\sigma_T^2+(\mu_T-\mu_S)^2}{2\sigma_S^2}-\frac12.
 \label{eq:kd-gaussian}
\end{equation}
这可以通过高斯对数密度相减后，对教师分布取期望得到，所用二阶矩为 $\mathbb E_T[(Y-\mu_S)^2]=\sigma_T^2+(\mu_T-\mu_S)^2$。若方差固定且相同，目标对均值的依赖退化为加权平方误差；若方差可学习，则还在匹配分布宽度，不能只解释成均值回归。

\keyterm{多标签分类（multi-label classification）}允许一个输入同时属于多个类别，例如一张图同时包含人和车。各类别可用 Bernoulli 概率建模，而不是用总和为一的 softmax 强制竞争。对单个类别，软监督是 $-q\log p-(1-q)\log(1-p)$，蒸馏可以对类别求和或平均。若采用独立 Bernoulli 假设，就没有显式描述类别之间的联合依赖；想迁移共现结构还需另外的关系建模。多类别与多标签仅一字之差，却对应不同概率空间。

\subsection{检测、分割与结构对应}
语义分割给每个像素预测类别，可将分类 KL 推广到有效像素集合，但需要处理忽略标签、边界区域、分辨率和类别不平衡。大面积背景容易主导平均损失，使小目标区域的监督被淹没。使用前景权重或区域采样会改变目标分布，选择这些设计应依据任务需求和验证结果，而非认为逐像素 KL 已经充分刻画所有结构。

目标检测还包含对象存在性、类别、边界框与多尺度位置。师生候选框数量或生成机制可能不同，必须先建立候选对应关系，再比较分类和位置；把两个未对齐的数组按同一索引相减没有合理语义。边界框可能用角点、中心和宽高等不同参数表示，坐标尺度和归一化也需一致。输出为集合时，还要区分集合元素的排列无关性与数组存储顺序。特征蒸馏可绕开部分输出对齐问题，却仍需要处理尺度和空间对应。

\subsection{Transformer 的隐藏状态和注意力}
单个注意力头可写成 $\mathbf Q=\mathbf H\mathbf W_Q$、$\mathbf K=\mathbf H\mathbf W_K$、$\mathbf V=\mathbf H\mathbf W_V$，其中 $\mathbf H\in\mathbb R^{L\times d_h}$，$\mathbf Q,\mathbf K\in\mathbb R^{L\times d_k}$，$\mathbf V\in\mathbb R^{L\times d_v}$。注意力矩阵为 $\mathbf A=\operatorname{softmax}_{\mathrm{row}}(\mathbf Q\mathbf K^\top/\sqrt{d_k}+\mathbf M)$，$\mathbf M$ 是位置掩码，输出为 $\mathbf A\mathbf V$。行 softmax 在允许关注的位置上归一化，因果掩码禁止观察未来词元。

蒸馏可比较隐藏状态、注意力矩阵或投影后的值关系。注意力矩阵维度取决于序列长度而不是隐藏宽度，因而在相同词元对齐条件下便于比较；但教师头数与学生头数不同时仍需指定头的配对或聚合。仅匹配 $\mathbf A$ 不等于匹配最终计算，因为值向量 $\mathbf V$ 也影响输出。注意力图相同也不能证明两个模型具有同样的内部解释或推理策略。

\subsection{跨模态与特权信息的可恢复边界}
多模态教师可以同时观察图像 $X$ 和额外模态 $Z$，学生部署时可能只观察 $X$。训练阶段可用而部署阶段不可用的信息称为\keyterm{特权信息（privileged information）}。教师利用深度传感器或额外文本获得更好的标签分布，学生通过蒸馏学习其中能由自身输入预测的部分。这要求模态之间有足够关联，而不是所有额外信息都能压缩到另一模态中。

例如教师输出向量 $\mathbf Q=q_T(X,Z)$，学生只预测 $g(X)$。在二阶矩有限时，平方误差满足条件期望分解
\[
 \mathbb E\|g(X)-\mathbf Q\|_2^2
 =\mathbb E\|g(X)-\mathbb E[\mathbf Q\mid X]\|_2^2
 +\mathbb E\|\mathbf Q-\mathbb E[\mathbf Q\mid X]\|_2^2.
\]
交叉项为零，因为 $\mathbb E[\mathbf Q-\mathbb E[\mathbf Q\mid X]\mid X]=0$。在不受函数类限制时，最优学生只能输出 $\mathbb E[\mathbf Q\mid X]$；第二项是仅观察 $X$ 时不可消除的残差。如果额外模态含有与 $X$ 独立且决定答案的信息，蒸馏不能让学生在每个样本上恢复它。这是信息可用性的限制，而不是训练技巧不足。

\subsection{扩散模型中的采样步骤蒸馏}
\keyterm{扩散模型（diffusion model）}通过噪声层级上的预测与迭代更新生成样本。蒸馏可以压缩网络，也可以在保持网络大小近似不变时减少采样步骤。若教师的确定性采样更新把状态 $x_t$ 先映到 $x_u$，再映到 $x_s$，学生可以学习一个从 $x_t$ 直接到 $x_s$ 的映射。概念性目标为
\[
 \mathbb E_{t,u,s,x_t}\left\|
 G_\theta(x_t,t,s)-\operatorname{sg}
 \bigl(F_\phi(F_\phi(x_t,t,u),u,s)\bigr)\right\|_2^2,
 \qquad t>u>s.
\]
这里 $F_\phi$ 表示指定教师采样器的一步状态更新，$G_\theta$ 是学生跨步映射；这只是函数复合层面的解释，不是所有扩散蒸馏方法共同使用的精确参数化。渐进蒸馏研究了逐阶段减少采样步数的方式\cite{salimans2022}。

这种场景的监督对象是动力学轨迹或去噪映射，分类 softmax 的温度解释不能直接照搬。训练需要覆盖学生将遇到的噪声状态，并关注多步误差累积。一次生成调用次数减少也不自动保证图像质量、多样性或条件控制保持不变，需要在固定采样设置下评价。教师采样步数、引导强度、噪声日程和学生更新形式都会影响被迁移的行为。

\section{训练实现：从公式到可复核的流程}
\label{sec:kd-implementation}

\subsection{最小离线分类蒸馏算法}
首先规定输入为训练集 $\mathcal D_\ell$、独立验证集、冻结教师 $f_\phi$、待训练学生 $f_\theta$、温度 $\tau$、权重 $\alpha$、学习率安排及最大训练步数。输出为按验证准则选择的学生参数。一个不依赖特定程序库的流程如下。
\begin{enumerate}
 \item 初始化学生和优化器，检查师生类别顺序、输入归一化与输出维度一致；用少量样本核对教师预测确实合理。
 \item 抽取小批量并应用已定义的数据增强。将教师置于约定的推理行为下，关闭其参数梯度，计算教师分数及软目标。
 \item 学生以前向训练行为计算分数，分别获得温度一的硬标签交叉熵与温度 $\tau$ 的师生 KL。先对类别求和，再对样本平均。
 \item 按式\eqref{eq:kd-classical}合并损失，清除旧梯度，对学生和训练所需适配器反向传播，再由优化器更新；教师保持固定。
 \item 定期在验证集上记录任务指标、概率指标及各损失分量，按预先规定的停止或选择规则保存学生。最终仅对选定模型运行测试评估。
\end{enumerate}
若使用梯度累积，每个微批量的缩放应保证最终等价于希望优化的样本平均；序列长度不同或最后一个微批量较小时，简单按微批次数平均未必正确。训练达到规定步数只意味着算法终止，不等于已经求得目标的全局最优解。

\subsection{损失缩放、精度与数值稳定性}
分类批量 KL 应明确写成 $B^{-1}\sum_{i=1}^B\sum_{k=1}^Kq_{ik}(\log q_{ik}-\log p_{ik})$。若程序对所有 $BK$ 个元素平均，就额外除以 $K$，软监督会随着类别数增加而减弱。语言模型中还要明确对词表先求和、对有效位置再平均。检查程序库的 reduction 语义比照抄某个权重数字更重要；库函数名称相似也不意味着默认约定相同。

使用稳定的对数 softmax 计算 $\log p$，避免先得到极小概率再取对数造成下溢。低精度训练时，概率归一化、KL 累加或高温下的微小分布差可以用更高精度计算。理论 KL 非负，但有限精度求和可能出现极小负值；较大的负值则往往意味着参数顺序、对数输入或归一化有误。简单把所有负值截成零会掩盖这些问题，也可能改变梯度。

若教师返回的是精确完整的温度一分布，且所有分量正，可以用 $q_k^\tau=(q_k^1)^{1/\tau}/\sum_j(q_j^1)^{1/\tau}$ 恢复指定温度的分布，因为未知 logit 平移会在归一化中抵消。但量化、截断或已经变成零的概率无法用此公式恢复丢失的信息。只存教师最大类别更不可能支持任意温度重新蒸馏，这应在缓存策略设计时提前考虑。

\subsection{截断概率与尾部质量}
设教师仅提供候选集合 $A$ 中的精确概率，尾部总质量为 $q_{\mathrm{tail}}=1-\sum_{k\in A}q_k$。学生仍可计算 $p_{\mathrm{tail}}=1-\sum_{k\in A}p_k$，从而构造“已知候选加其他”这个较粗类别空间上的 KL。由第\ref{sec:kd-response}节同样的分组分解，完整 KL 等于粗粒度 KL 加 $q_{\mathrm{tail}}$ 乘以尾部条件 KL，因此粗粒度 KL 是完整 KL 的下界。

这个构造保留了尾部总质量，却没有区分尾部内部的类别。直接把保留候选重新归一化，相当于条件化到候选集合内，又定义了不同目标。若教师只返回已经重新归一化的截断概率，连原始尾部总质量也无法精确知道。报告使用 top-$k$ 概率蒸馏时，应说明候选如何选取、是否包含学生高概率项、尾部怎样处理以及监督损失属于哪个事件空间。

\subsection{梯度冲突与训练日程}
记硬标签梯度为 $\mathbf g_h$、软标签梯度为 $\mathbf g_s$。若两者内积为负，沿其中一个负梯度的很小更新可能增加另一项损失；这是多目标优化中的局部冲突。梯度范数和夹角可以帮助诊断软项是否压倒任务项，但不能仅凭一次内积为负就删除蒸馏信号，因为教师也可能纠正噪声标签。应结合样本内容、验证趋势和长期训练行为判断。

逐步增加蒸馏权重、先拟合中间提示再联合训练、先监督预热再蒸馏，都是可能的训练日程。它们改变优化路径，需要与相同预算的直接联合训练比较。教师过强而学生过小时，可尝试中等容量的中间教师，但额外阶段会增加成本，且不保证逐级逼近总比直接逼近好。合理的做法是让这些设计回应已观察到的优化或容量问题，而非把所有技巧一次叠加后无法解释收益来源。

\section{实验设计、评价指标与成本分析}
\label{sec:kd-evaluation}

\subsection{首先建立什么对照}
一项蒸馏实验至少需要知道教师性能、相同学生仅用真实标签训练的性能，以及学生蒸馏后的性能。若没有无蒸馏学生基线，无法判断提升来自蒸馏还是学生本身的架构与训练改进。进一步可设置标签平滑、教师硬伪标签、完整软目标和特征辅助等对照，每次只改变一个主要因素。数据增强、初始化、训练步数、学习率搜索和验证预算应尽量匹配，或明确列出无法匹配的差异。

对生成任务，教师文本带来的额外数据量可能本身就有收益。需要比较相同数量的原始任务数据、教师生成数据，以及不同筛选策略；最好同时控制样本数和有效词元数。若一个方法生成多条答案后选最好的一条，而基线只生成一次，性能差异包含额外推理搜索的贡献。评测提示、解码参数、最大输出长度、停止规则和答案抽取方法都应保持可比。

消融实验回答“去掉某个组件后发生什么”，并不自动给出唯一因果机制。某组件移除后，如果其他权重不重新选择，可能只是整个目标的数值比例失衡。应先说明比较的是固定超参数下的组件作用，还是在相同调参预算下各方案的最佳表现。两者都可以有价值，但对应不同问题。训练失败或性能下降的设置也应记录，避免只展示筛选后的成功点。

\subsection{准确率之外还应看什么}
分类准确率为 $n^{-1}\sum_i\mathbf1\{\arg\max_kp_k(x_i)=y_i\}$；负对数似然为 $-n^{-1}\sum_i\log p_{y_i}(x_i)$。前者对概率大小不敏感，后者会惩罚自信错误。还可报告 Brier 分数 $n^{-1}\sum_i\sum_k(p_k(x_i)-\mathbf1\{y_i=k\})^2$，以平方形式评价概率预测。蒸馏可能改善其中一个指标而恶化另一个，所以应依据使用目标预先确定主指标。

常见的\keyterm{期望校准误差（expected calibration error, ECE）}把最大预测概率划分为若干区间，对每个非空区间 $I_b$ 计算准确率与平均置信度之差，再求加权绝对值：
\[
 \widehat{\operatorname{ECE}}=\sum_b\frac{|I_b|}{n}
 \left|\operatorname{acc}(I_b)-\operatorname{conf}(I_b)\right|.
\]
这是依赖分箱的样本指标，结果会受区间数量、样本量和分箱方式影响；它也只观察最大类别置信度这一种校准性质。不能把一个很小的 ECE 当成所有类别、所有子群都已经校准的证明。

还应按类别、难度、长度或场景报告分组结果，尤其观察教师出错而学生是否被带错的样本。分布外测试需要明确偏移方式，例如新域、噪声、遮挡或组合变化；不能用一个受损图像集概括所有鲁棒性。对于不均衡类别，宏平均指标与总体样本平均有不同含义，应按照任务要求选择。模型与教师更一致可能同时使真实标签准确率下降，因此一致率只能作为过程指标。

\subsection{随机性、置信区间与数据泄漏}
神经网络训练受初始化、批次顺序和随机增强影响。重复多个随机种子可帮助评估稳定性，但需要同时报告次数、均值和离散程度。只有一次小幅提升时，证据通常不足以判断方法稳定优于基线。对相同测试样本比较两个模型，应利用配对结果，而不是把两组误差当成来自独立样本；对具有用户或文档簇的数据，还应按独立单位划分和重采样。

蒸馏引入教师这一额外信息来源，也增加了数据泄漏的路径。教师训练集可能包含测试样本，合成提示可能直接复制测试题，缓存目标可能使用了不可用于训练的标签。实验应明确教师数据已知到什么程度，并区分确定的无泄漏协议与无法完全确认的条件。这里的重点是评价边界：若教师训练数据不透明，就不应把测试表现解释为已经严格证明对完全新知识的泛化。

\subsection{参数量、延迟与总计算成本}
设教师单次预测成本为 $c_T$，学生为 $c_S$，额外蒸馏训练和目标生成成本为 $C_{\mathrm{extra}}$，预期部署调用次数为 $N$。若其他条件相同且 $c_T>c_S$，简化成本模型要求
\begin{equation}
 N(c_T-c_S)>C_{\mathrm{extra}}
 \label{eq:kd-amortization}
\end{equation}
才会在累计成本上抵消额外训练开销。若教师已存在，是否计入其原始训练成本取决于比较问题；如果比较从零构建两套系统，就应把它纳入。上述是统一计量单位下的成本关系，不预测具体硬件价格或真实账单。

实际速度应在相同设备、批量、精度和输入输出长度下测量，并区分首个输出延迟、后续生成吞吐量和完整请求延迟。冷启动、预热和并发也会影响结果。语言模型的键值缓存消耗、蒸馏后回答长度变化，以及检测模型的后处理都可能主导端到端表现。参数量、理论浮点运算次数和实测延迟应分别报告，不能用其中一个完全代替另外两个。

\subsection{与量化、剪枝组合时的比较}
\keyterm{量化（quantization）}将权重或激活映射到较少数值等级，\keyterm{剪枝（pruning）}则删除参数或结构。蒸馏可以为受这些约束的学生提供训练信号，使其在压缩后重新拟合任务行为。但训练时的近似计算与实际部署算子可能不同，验证必须针对最终导出的模型。量化误差、结构改变和蒸馏监督是不同因素，应分别设置对照，避免把三者合用的结果全归因于其中一个。

若采用非结构化稀疏，权重零元素很多也不意味着普通密集算子自动加速；若采用低比特量化，也需设备支持相应算子。蒸馏不能替代这些系统条件。最终有意义的结论应形如“在规定精度与延迟约束下，某个学生配置达到怎样的任务表现”，并说明训练开销和评估范围，而不是只说知识被压缩了多少倍。

\section{贯穿算例与完整实验方案}
\label{sec:kd-example-project}

\subsection{手算一次混合监督更新}
继续使用 $\mathbf q=(0.7,0.2,0.1)$、$\mathbf p=(0.6,0.3,0.1)$、真实类别一，取 $\tau=1$、$\alpha=0.4$。有效目标由式\eqref{eq:kd-mixed-target}给出，等于 $(0.88,0.08,0.04)$，所以混合损失对三个 logit 的梯度为 $(-0.28,0.22,0.06)$。纯硬标签梯度则为 $(-0.4,0.3,0.1)$。软监督没有要求完全放弃第二类和第三类，而是让这两个方向的推低幅度减小。

取初始分数 $\mathbf z=(\log0.6,\log0.3,\log0.1)$，softmax 恰好还原 $\mathbf p$。若直接把分数当作参数，以学习率 $\eta=0.1$ 做一步梯度下降，则新分数为 $\mathbf z'=(\log0.6+0.028,\log0.3-0.022,\log0.1-0.006)$。再做 softmax，得到新概率约为 $(0.610982,0.290592,0.098426)$，正确类别概率有所提高。这里不能直接把概率减去梯度，否则可能破坏非负性和归一化。在完整网络中，logit 还依赖共享参数，实际改变由网络 Jacobian 决定，不能把这个逐点计算当成整个模型的一次精确更新。

当学生恰好等于有效目标时，此样本的输出梯度为零；当学生恰好等于教师时，硬标签项仍可能产生非零梯度。例如 $\mathbf p=\mathbf q$ 时，混合梯度是 $0.6(\mathbf q-\mathbf e_1)$。这解释了为何同时使用真实标签后，训练目标通常不是把学生固定在教师分布上。教师一致性与任务监督共同决定折中位置。

\subsection{一个可复现的小型分类研究}
可用二维三类合成数据建立最小实验：三个类别分别从具有重叠区域的高斯分布抽样，明确记录各类均值、协方差、先验和随机种子。由于生成分布已知，可以计算真实条件类别概率，因而能分别观察教师估计误差、学生模仿误差和对真实分布的误差。训练、验证和测试样本独立生成，避免仅在训练点可视化决策边界而误判泛化。

教师选择较宽的多层感知机，学生选择较窄网络，再加入线性学生检验容量限制。先训练教师和无蒸馏学生，确认优化设置合理；随后固定教师，仅改变温度和软监督权重。可以在温度集合 $\{1,2,4,8\}$、权重集合 $\{0,0.25,0.5,0.75\}$ 上用验证集比较，选择规则预先固定。这里的集合只是覆盖不同强度的示例搜索空间，并非对任何任务都推荐的最佳范围。

记录测试准确率、负对数似然、真实条件分布 KL、教师条件分布 KL，以及训练时间。将蒸馏数据局限到平面的一侧，可检验覆盖不足；人为扰动教师某个区域的标签，可检验错误传递；换成理想条件分布教师，可分离教师误差与学生容量问题。这些干预让理论概念对应到可观察现象，比只在一个大型数据集上堆叠复杂组件更适合建立理解。本文给出方案，不预设各组实验的胜负。

\subsection{故障诊断的因果顺序}
如果蒸馏后学生退化，先核对类别顺序、输入预处理、温度、损失归约和教师冻结；这些错误会使一个合理方法在实现上失真。再看硬、软损失的量级和梯度，确认教师目标确实在影响训练。随后按教师正确与错误的样本分组，判断退化来自监督偏差还是整体优化不稳定。最后再尝试调整温度、权重、数据覆盖或蒸馏载体，避免在基础实现错误上反复调参。

如果软 KL 已经很低而任务仍差，应考虑教师本身不可靠、训练分布与测试分布不同、类别间隔很小或评测协议错误。如果训练 KL 始终很高，应考虑学生容量不足、优化未收敛、目标缓存错位或特征约束过强。若训练很慢，应分别测量教师前向、学生前后向、目标存取和数据加载，先确定瓶颈再决定缓存或压缩目标。诊断依据是可测量现象与具体机制之间的对应，而不是为每次失败换一个新方法名称。

\section{练习与参考解答}
\label{sec:kd-exercises}

\subsection{练习一：概率相同而分数不同}
设教师分数为 $(3,1,-1)$，学生分数为 $(8,6,4)$。证明任意正温度下两者概率相同，并解释直接 logit 均方误差为什么不能识别这种等价。

\begin{solve}
学生分数等于教师分数加 $5\mathbf1$，共同因子 $e^{5/\tau}$ 在 softmax 分子分母中抵消，所以两分布一致，KL 为零。三个坐标差均为五，按坐标平均的平方误差为二十五。分别减去各自均值后，二者都等于 $(2,0,-2)$，中心化平方误差才为零。由此可见，概率目标约束相对分数，未中心化目标还额外约束不可由概率识别的共同偏移。
\end{solve}

\subsection{练习二：教师错误的权重阈值}
在二分类问题中，真实标签为第一类，教师输出 $(0.2,0.8)$，温度为一。求混合有效目标仍把第一类作为唯一最大类别时，软监督权重 $\alpha$ 的范围。

\begin{solve}
有效目标为 $(1-0.8\alpha,0.8\alpha)$。第一分量严格大于第二分量当且仅当 $1>1.6\alpha$，所以 $0\leq\alpha<0.625$。等于该值时两类打平，超过时教师错误类别成为最大分量。阈值源于自由输出的逐点最优目标；共享参数网络是否实际达到它，还取决于其他样本和优化。
\end{solve}

\subsection{练习三：平方补偿是否使梯度完全不变}
有人认为蒸馏损失乘以 $\tau^2$ 后，改变温度不再影响训练。指出其推理错误，并给出需要检查的量。

\begin{solve}
补偿后的梯度是 $\tau(\mathbf p^\tau-\mathbf q^\tau)$，其中两个分布本身仍依赖温度。只有在固定有限分数的高温极限中，它才趋于 $K^{-1}(\widetilde{\mathbf z}-\widetilde{\mathbf a})$。有限温度改变各类别相对权重，训练过程又会改变分数尺度，因此不存在一般的梯度完全不变结论。应检查中心化分数相对温度的尺度，以及不同温度下软梯度与硬梯度的实际范数。
\end{solve}

\subsection{练习四：合并尾部为什么低估完整 KL}
设类别被分为已知候选和尾部，证明只比较尾部总质量的 KL 不大于完整 KL，并说明何时相等。

\begin{solve}
若尾部质量 $q_A,p_A$ 均正，把尾部各项写成 $q_k=q_A\bar q_k$、$p_k=p_A\bar p_k$，其贡献为 $q_A\log(q_A/p_A)+q_AD_{\mathrm{KL}}(\bar{\mathbf q}\|\bar{\mathbf p})$。粗粒度损失只保留第一项，漏掉的第二项非负。若教师尾部质量正且散度有限，等号要求师生尾部条件分布相同；教师尾部质量为零时，尾部内部无需额外匹配。学生尾部质量为零而教师为正时，两种散度均按约定处理为无穷，不能做普通有限差值计算。
\end{solve}

\subsection{练习五：参考前缀与学生前缀}
一个学生在人工参考前缀上与教师的词元 KL 很低，但自由生成经常失败。给出两种机制解释，并说明在策略蒸馏能直接针对哪一部分。

\begin{solve}
第一种解释是学生早期偶然生成不同词元，进入参考数据没有覆盖的前缀，后续条件分布便缺少监督；在学生生成前缀上查询教师直接针对这种状态覆盖差异。第二种解释是教师本身在这些任务上给出错误分布，或者学生解码与评价规则不合适；只改变前缀来源不能自动解决教师错误和解码协议问题。因此应同时观察前缀分布、教师反馈质量和生成评价，而非把所有失败统称为暴露偏差。
\end{solve}

\subsection{练习六：额外模态是否总能蒸馏进去}
设 $X$ 与二元随机变量 $Z$ 独立，$P(Z=0)=P(Z=1)=1/2$，标签 $Y=Z$。教师同时观察 $X,Z$，学生仅观察 $X$。学生能否在新样本上达到教师的完美准确率？

\begin{solve}
不能。独立性给出 $P(Y=1\mid X)=1/2$，任何仅依赖 $X$ 的确定分类器在该分布上的准确率都是 $1/2$，独立随机化也不能提高其期望准确率。教师因直接观察 $Z$ 可以总是正确。蒸馏增加了训练监督，却没有给部署学生增加关于当前新样本 $Z$ 的观测，因而无法消除这个信息差距。这不否定跨模态蒸馏的价值，只说明它依赖可由学生输入预测的关联结构。
\end{solve}

\section{阅读原始文献的方法与知识边界}
\label{sec:kd-reading}
阅读一篇蒸馏论文，可以先把其方法写成四个要素：用于查询教师的输入分布、教师暴露的知识载体、师生之间的比较函数、教师与学生的更新关系。经典输出蒸馏主要改变概率监督，特征方法改变比较对象，关系方法改变对样本集合的约束，在线与自蒸馏改变目标来源，在策略方法改变被访问的前缀分布。不同方向可以组合，名称不同也可能只改变其中一个要素；这种分解比仅记缩写更有利于比较机制。

其次，应把数学恒等式与经验主张分开。KL 分解、高温极限和条件期望分解可以在明确条件下证明；某类中间表示更适合迁移、某教师更适合某学生、某个温度带来最佳性能，则需要特定任务实验。学生超越教师、无标签训练有效或采样加速等结论，必须连同数据、结构、预算和评估协议一起理解。本文讨论的是有代表性的基础机制和扩展方式，不声称穷尽所有任务中的蒸馏算法。

文献\cite{gou2021}可用于了解研究版图；入门推导可先对照文献\cite{hinton2015}，再分别阅读特征、关系和训练组织方式的原始论文。进入语言模型时，先读序列级蒸馏\cite{kim2016}并确认生成数据的分布，再读在策略蒸馏\cite{agarwal2024}以理解训练状态分布的变化。阅读实验表格时，同时寻找失败条件、计算预算和消融设置；阅读公式时，核对每个期望的采样分布、每个梯度的变量以及哪些分支被视为常量。能够把这些条件写清楚，才算把蒸馏原理转化为可检验的训练方案。

\begin{thebibliography}{99}
\bibitem{hinton2015}
Geoffrey Hinton, Oriol Vinyals, Jeff Dean.
Distilling the Knowledge in a Neural Network.
2015. \url{https://arxiv.org/abs/1503.02531}.

\bibitem{gou2021}
Jianping Gou, Baosheng Yu, Stephen John Maybank, Dacheng Tao.
Knowledge Distillation: A Survey.
International Journal of Computer Vision, 2021.
\url{https://arxiv.org/abs/2006.05525}.

\bibitem{romero2015}
Adriana Romero, Nicolas Ballas, Samira Ebrahimi Kahou, Antoine Chassang, Carlo Gatta, Yoshua Bengio.
FitNets: Hints for Thin Deep Nets.
ICLR, 2015. \url{https://arxiv.org/abs/1412.6550}.

\bibitem{zagoruyko2017}
Sergey Zagoruyko, Nikos Komodakis.
Paying More Attention to Attention: Improving the Performance of Convolutional Neural Networks via Attention Transfer.
ICLR, 2017. \url{https://arxiv.org/abs/1612.03928}.

\bibitem{park2019}
Wonpyo Park, Dongju Kim, Yan Lu, Minsu Cho.
Relational Knowledge Distillation.
CVPR, 2019. \url{https://arxiv.org/abs/1904.05068}.

\bibitem{tian2020}
Yonglong Tian, Dilip Krishnan, Phillip Isola.
Contrastive Representation Distillation.
ICLR, 2020. \url{https://arxiv.org/abs/1910.10699}.

\bibitem{zhao2022}
Borui Zhao, Quan Cui, Renjie Song, Yiyu Qiu, Jiajun Liang.
Decoupled Knowledge Distillation.
CVPR, 2022. \url{https://arxiv.org/abs/2203.08679}.

\bibitem{zhang2018}
Ying Zhang, Tao Xiang, Timothy M. Hospedales, Huchuan Lu.
Deep Mutual Learning.
CVPR, 2018. \url{https://arxiv.org/abs/1706.00384}.

\bibitem{furlanello2018}
Tommaso Furlanello, Zachary C. Lipton, Michael Tschannen, Laurent Itti, Anima Anandkumar.
Born Again Neural Networks.
ICML, 2018. \url{https://arxiv.org/abs/1805.04770}.

\bibitem{tarvainen2017}
Antti Tarvainen, Harri Valpola.
Mean Teachers Are Better Role Models: Weight-Averaged Consistency Targets Improve Semi-Supervised Deep Learning Results.
NeurIPS, 2017. \url{https://arxiv.org/abs/1703.01780}.

\bibitem{caron2021}
Mathilde Caron 等.
Emerging Properties in Self-Supervised Vision Transformers.
ICCV, 2021. \url{https://arxiv.org/abs/2104.14294}.

\bibitem{yin2020}
Hongxu Yin 等.
Dreaming to Distill: Data-Free Knowledge Transfer via DeepInversion.
CVPR, 2020. \url{https://arxiv.org/abs/1912.08795}.

\bibitem{kim2016}
Yoon Kim, Alexander M. Rush.
Sequence-Level Knowledge Distillation.
EMNLP, 2016, pp.~1317--1327.
\url{https://aclanthology.org/D16-1139/}.

\bibitem{agarwal2024}
Rishabh Agarwal, Nino Vieillard, Yongchao Zhou, Piotr Stanczyk, Sabela Ramos, Matthieu Geist, Olivier Bachem.
On-Policy Distillation of Language Models: Learning from Self-Generated Mistakes.
ICLR, 2024. \url{https://arxiv.org/abs/2306.13649}.

\bibitem{salimans2022}
Tim Salimans, Jonathan Ho.
Progressive Distillation for Fast Sampling of Diffusion Models.
ICLR, 2022. \url{https://arxiv.org/abs/2202.00512}.
\end{thebibliography}

\end{document}
